\documentclass{report}
\usepackage{amsfonts, amsmath}
\newcommand\R{{\mathbb R}}
\newcommand\Q{{\mathbb Q}}
\newcommand\N{{\mathbb N}}
\newcommand{\vf}{\varphi}
\pagestyle{empty}
\begin{document}
\hfuzz 12pt
\large
\centerline{\Huge \bf Analisi Matematica II}
\vskip.2cm
\centerline{\Huge Correzione Primo Appello}

\vskip1cm
\begin{enumerate}
\item Si calcoli il rapporto incrementale
\[
\lim_{h\to 0}\frac{f(x+h)-f(x)}h=\lim_{h\to
0}\int_0^1\frac{e^{-hy}-1}he^{-xy}\sqrt{1-y^2}dy.
\]
Ora \`e facile controllare che la succesione di funzioni
$\{\frac{e^{hy}-1}h\}$ converge uniformemente per $y\in[0,1]$, dunque
si puo portare il limite sotto il segno di integrale ottenendo
\[
f'(x)=-\int_0^1ye^{-xy}\sqrt{1-y^2}dy.
\]
\item Cominciamo con lo scrivere
\[
\int_{-\frac\pi 2}^{\frac\pi 2}e^{\sin x} dx=\sum_{n=0}^\infty\frac
1{n!}\int_{-\frac\pi 2}^{\frac\pi 2} (\sin x)^n dx
\]
dove abbiamo usato la uniforme convergenza della serie per portarla
fuori dal segno di integrale. Ora si noti che l'integrale da zero per
tutte le potenze dispari. Inoltre, per $n$ pari,
\[
\int_{-\frac\pi 2}^{\frac\pi 2} (\sin x)^n dx\leq \pi.
\]
Poich\`e
\[
\pi\sum_{n=4}^\infty \frac 1{n!}=\pi(e-2-\frac 12-\frac1 6)\leq .2
\]
segue che basta calcolare i primi tre termini della serie per avere un
risultato con la precisione desiderata. Cosi facendo si ottiene
\[
\pi+\frac 12\int_{-\frac\pi 2}^{\frac\pi 2} (\sin x)^2 dx=\pi+\frac\pi 4.
\]
\item Usando il criterio integrale si ottiene
\[
\sum_{n=2}^\infty\frac1{n(\ln n)^2}\leq \frac 1{2(\ln
2)^2}+\lim_{n\to\infty}\int_{2}^n \frac{1}{x(\ln x)^2} dx.
\]
Ora si operi la sostituzione $y=\ln x$, dunque
\[
\sum_{n=2}^\infty\frac1{n(\ln n)^2}\leq \frac 1{2(\ln
2)^2}+\lim_{n\to\infty}\int_{\ln 2}^{\ln n} \frac{1}{y^2} dy\leq 
\frac 1{2(\ln 2)^2}+\frac 1{\ln 2},
\]
dunque la serie convege.
\item Si provi con una soluzione del tipo $x(t)=e^{\alpha t}$
sostituendo si trova che la cosa funziona se $\alpha^2=2-\alpha$, cioe
ci sono due soluzioni $\alpha_+=\frac{-1+\sqrt 5}{2}$ e
$\alpha_-=\frac{-1-\sqrt 5}{2}$.
La soluzione generale sara allora della forma
$x(t)=ae^{\alpha_+t}+be^{\alpha_-t}$. Le costanti $a$ e $b$ devono
essere determinate usando le condizioni iniziali: $x(0)=0$ implica
$a=-b$, mentre $x'(0)=1$ implica $a=\frac 1{\sqrt 5}$.
\item Questa siottine facilmente utilizzando la formula per la
soluzione generale di queste equzioni. Oppure notando che 
\[
\frac{\sin t-\cos t}2
\]
\`e una soluzione della equazione differenziale. Questo significa che
la soluzione generale ha la forma
\[
x(t)=ae^t+\frac{\sin t-\cos t}2
\]
e aggiustando $a$ in modo da soddisfare la condizione iniziale.
\item Ovviamente per $|x|>1$ la serie non converge, visto
che il termine $n$-esimo non converge a zero. Per $|x|<1$ invece la
serie \`e maggiorata da una serie geometrica di ragione $|x|$ e dunque
converge. Il raggio di convergenza \`e dunqe $1$.
\end{enumerate}
\end{document}
